\documentclass[10pt,a4paper]{amsart}
\usepackage[margin=19mm]{geometry}
\usepackage{amsmath,amssymb,mathtools}
\usepackage[scr=rsfs,cal=euler]{mathalfa}
\usepackage{xfrac}
\usepackage[unicode,hidelinks,bookmarks=false]{hyperref}
\newcommand{\ie}{i.e.\ }
\newcommand{\cf}{cf.\ }
\newcommand{\resp}{respectively\ }
\newcommand{\Subsection}{\subsection}
\providecommand{\nobreakdash}{\nobreak\hbox{-}}
\setlength{\emergencystretch}{2em}
\multlinegap0pt
\usepackage[shortlabels]{enumitem}
\usepackage{xcolor}
\usepackage{amssymb}
\usepackage{mathtools}
\newtheorem{theorem}{Theorem}
\newcommand\R{\mathbb{R}}
\renewcommand{\ge}{\geqslant}
\renewcommand{\geq}{\geqslant}
\renewcommand{\le}{\leqslant}
\renewcommand{\leq}{\leqslant}
\newcommand\ve{\varepsilon}
\title{Sums of three powerful numbers}
\author{Tim Browning, Matteo Verzobio}
\date{Source: arXiv:2608.24512v1 (CC BY 4.0)}
\begin{document}
\maketitle

\noindent\emph{Source and licence.} Sums of three powerful numbers. Version arXiv:2608.24512v1 (CC BY 4.0), distributed by arXiv under the Creative Commons Attribution 4.0 International licence, http://creativecommons.org/licenses/by/4.0/, as recorded in the arXiv licence metadata for this exact version and shown on the abstract page https://arxiv.org/abs/2608.24512v1. Full licence notice: \texttt{licenses/arxiv-2608.24512-CC-BY-4.0.txt}.\par\smallskip

\noindent\textit{Editorial scope.} Counted unit: the article's theorem prop:etagen (source0.tex lines 1026--1054). The article's complete original proof of it is delivered with it (source0.tex lines 1056--1190, one \texttt{proof} environment of 135 lines). No mathematics is written, repaired or re-derived: the statement and the proof are the article's own lines, in the article's own notation, and no retained line is substituted or re-flowed.  Retained uncounted as the source-local context the proof needs: lines 167--175; lines 241--243.  Nothing was omitted from any retained span, so the package carries no omission record.  No retained line contains a citation, so the wrapper carries no bibliography and no bibliography entry.  No retained line contains a figure environment, an image-inclusion command or drawing code, so no image is delivered and the package declares no asset dependency.  Editorial formatting only, in addition: an AMS \texttt{amsart} preamble that restates the article's own declarations and macros used by the retained lines, namely the environments \texttt{theorem} on separate counters, together with the standard packages those lines need.  The retained environments print in this document as Theorem 1 in order of appearance, whereas the article prints the same items with the numbering of its own full document.  The complete native source of the article as distributed in the arXiv e-print for this exact version is bundled unchanged as \texttt{sources/208484/source0.tex}.
\begin{equation}\label{eq:hawk}
N(B)
 =
 \#\left\{
 (a,b,c)\in
 (\mathcal S_p\times\mathcal S_q\times\mathcal S_r)\cap[1,B]^3:
 \gcd(a,b,c)=1,\ a+b=c
 \right\}.
\end{equation}

\begin{equation}\label{eq:triv}
N(B)\ll_{p,q}B^{1/p+1/q},
\end{equation}

\begin{theorem}\label{prop:etagen}
Let $p\geq q\geq r\ge 2$. Assume that there is a fixed $\kappa\geq 0$ such that, uniformly for nonzero $\lambda,\mu,\nu\in\mathbb Z$, for every $0\le k\le r-1$, for all $X,Y,Z\geq 2$ and $\ve>0$, we have 
\[
\#\left\{(x,y,z)\in \mathbb N^3:
\begin{array}{l}
x\le X,\ y\le Y,\ z\le Z,\ \gcd(x,y,z)=1,\\
\lambda x^p+\mu y^q=\nu z^{r+k}
\end{array}
\right\}
\ll_{\varepsilon,p}
H^{\ve}\max\{X,Y\}^{\kappa}
\]
with $H=\max\{X,Y,Z\}$.
Then 
\[
N(B)\ll_{\delta, p}B^{\frac{1}{p}+\frac{1}{q}-\delta},
\]
for any 
\[
\delta<
 \frac{
 \displaystyle
 \frac1p+\frac1q-
 \left(
 \frac1r-\frac1{qr}+\frac{\kappa}q
 \right)}
 {p+q+1+1/r}.
\] 
\end{theorem}

\begin{proof}
Any nonzero $m$-full integer $n\in \mathbb{N}$ can be written uniquely in the form
\[
n = v_0^{m} \prod_{s = 1}^{m-1} v_s^{m + s},
\]
for $v_0, v_1, \ldots, v_{m - 1} \in \mathbb{N}$, such that $\mu^2(v_s) =1$ for $1 \leq s \leq m-1$ and $\gcd(v_s, v_{s'}) = 1$ for $1 \leq s < s' \leq m-1$.
Recall 
from~\eqref{eq:hawk} that 
\[
N(B)=\#\{(a,b,c)\in (\mathcal{S}_p\times \mathcal{S}_q\times \mathcal{S}_r)\cap[1,B]^3: \gcd(a,b,c)=1,~a+b=c\}.
\]
We adopt the factorisation
\[
a=x_0^p\prod_{i=1}^{p-1}x_i^{p+i},\qquad
b=y_0^q\prod_{j=1}^{q-1}y_j^{q+j},\qquad
c=z_0^r\prod_{k=1}^{r-1}z_k^{r+k}.
\]
The squarefree and coprimality conditions in this parametrisation make it unique;
for upper bounds, they can only reduce the number of admissible tuples. 

We decompose all variables dyadically. On a fixed dyadic box, write
\[
x_i\sim B^{\alpha_i},\qquad y_j\sim B^{\beta_j},\qquad z_k\sim B^{\gamma_k}.
\]
Let
\begin{equation}\label{eq:UVW}
S=\sum_{i=0}^{p-1}\alpha_i,\qquad
T=\sum_{j=0}^{q-1}\beta_j,\qquad
U=\sum_{k=0}^{r-1}\gamma_k.
\end{equation}
The conditions $a,b,c\le B$ imply
\[
\sum_{i=0}^{p-1}(p+i)\alpha_i\le 1,\qquad
\sum_{j=0}^{q-1}(q+j)\beta_j\le 1,\qquad
\sum_{k=0}^{r-1}(r+k)\gamma_k\le 1.
\]
In particular, we have 
\[
S\le \frac1p,\qquad T\le \frac1q,\qquad U\le \frac1r.
\]

Fix a dyadic box, and suppose that its contribution is
$O_{p}(B^{\Delta})$ on this box. 
The argument behind the trivial bound~\eqref{eq:triv} easily gives
\begin{equation}\label{eq:triv'}
\Delta\le S+T,\qquad \Delta\le S+U,\qquad \Delta\le T+U.
\end{equation}
The fourth bound we will use comes from fixing every variable except $x_0,y_0,z_k$, for a choice of $0\le k\le r-1$. The equation becomes
\[
\lambda x_0^p+\mu y_0^q=\nu z_k^{r+k}
\]
with fixed nonzero coefficients $\lambda,\mu,\nu$. Since $\gcd(a,b,c)=1$, we have
$\gcd(x_0,y_0,z_k)=1$. Therefore, for every $0\le k\le r-1$, it follows from the hypothesis of the theorem that 
\begin{equation}\label{eq:4way}
\Delta
\le
S+T+U-\alpha_0-\beta_0-\gamma_k
+\kappa\max\{\alpha_0,\beta_0\}+\ve.
\end{equation}

Suppose that
\[
\Delta> \frac 1p+\frac 1q-\eta,
\]
where 
\[
\eta=
 \frac{
 \displaystyle
 \frac1p+\frac1q-
 \left(
 \frac1r-\frac1{qr}+\frac{\kappa}q
 \right)}
 {p+q+1+1/r}\in \R. 
\]
Since $\Delta\le S+T$, by~\eqref{eq:triv'}, 
while $S\le 1/p$ and $T\le 1/q$, it follows that 
\[
S\ge \frac1p-\eta,\qquad T\ge \frac1q-\eta.
\]
Moreover, since $S=\alpha_0+\sum_{i=1}^{p-1}\alpha_i$, we have
$
1\ge pS+(S-\alpha_0).
$
Hence
\begin{equation}\label{eq:Ua0}
 S-\alpha_0\le 1-pS\le p\eta.
\end{equation}
Similarly,
\begin{equation}\label{eq:Vb0}
T-\beta_0\le q\eta.
\end{equation}

The bound $\Delta\le S+U$ 
in~\eqref{eq:triv'}
also gives
$
U\ge \Delta-S\ge \frac1q-\eta.
$
Let
$M=\max_{0\le k\le r-1}\gamma_k$,
and choose $k$ such that $\gamma_k=M$. Since there are $r$ variables on the $z$-side, we obtain
\begin{equation}\label{eq:MWr}
 M\ge \frac Ur\ge \frac1{qr}-\frac{\eta}{r}.
\end{equation}
Combining~\eqref{eq:4way},~\eqref{eq:Ua0} and~\eqref{eq:Vb0}, we obtain
\[
\Delta
\le
(p+q)\eta+(U-M)
+\kappa\max\{\alpha_0,\beta_0\}+\ve.
\]
Since $p\ge q$, we have 
$
\max\{\alpha_0,\beta_0\}
\le
\frac 1q.
$
Substituting this and~\eqref{eq:MWr} into the bound for $\Delta$, we obtain
\begin{equation}\label{eq:finaleta}
\Delta
\le
(p+q)\eta+
\frac1r-\frac1{qr}+\frac{\eta}{r}+\frac{\kappa}{q}+\ve=
\frac 1p+\frac 1q-\eta+\ve.
\end{equation}
Recall that we were assuming that $\Delta> \frac 1p+\frac 1q-\eta $. But if not, then~\eqref{eq:finaleta} clearly holds.
Thus, no dyadic box can contribute with exponent more than $\frac1p+\frac1q-\eta+\ve$.
The number of dyadic boxes is $O_{p}((\log B)^{p+q+r})$, which is
absorbed into $B^\varepsilon$. Hence
\[
N(B)\ll_{\varepsilon,p}B^{\frac 1p+\frac 1q-\eta+\varepsilon},
\]
for any $\ve>0$, which thereby completes the proof. 
\end{proof}

\end{document}
